\section{Scalar estimates and the characteristic identities}
\label{sec:analytic}

This section supplies two kinds of global information: analytic bounds
for the sizes of the Weyl blocks, and integer identities relating their
second moments to volume. The former will control derivatives in the
coupled inequality; the latter will restrict the possible imbalance
between the two blocks.

The standing setting is the closed, connected, oriented Einstein
four-manifold with $\operatorname{Ric}=3g$ fixed in the introduction.
Appendix~\ref{app:analytic} supplies the analytic foundations used here;
its proofs depend only on this setting and the curvature conventions of
Section~\ref{sec:preliminaries}. We retain the weighted variance argument
and the complete Weyl-gap proof in the main text because their precise
weights and behavior at zeros are used in the two-block argument.

\subsection{A strengthened scalar variance estimate}

\begin{proposition}[Weighted variance]
\label{prop:weighted-variance}
For every real $a\in H^1(M)$,
\begin{equation}
 \operatorname{Var}(a)
 :=\mathbf E\bigl[(a-\mathbf E a)^2\bigr]
 \le\mathbf E\frac{|\nabla a|^2}{K},
 \qquad K=4+v^2+b^2.
 \tag{4}\label{eq:4}
\end{equation}
\end{proposition}

\begin{proof}
First let $\phi$ be smooth and write
$S=\operatorname{Hess}_0\phi$. The integrated Bochner identity
and $\operatorname{Ric}=3g$ give
\[
 \mathbf E\bigl[(\Delta\phi)^2\bigr]
 =\mathbf E\left(4|\nabla\phi|^2+\frac43|S|^2\right).
\]
Define operators on covariant two-tensors by
\[
 (T_D)_{ikjl}=\frac12\sum_s t_s I_{s,ik}I_{s,jl},
 \qquad
 (U_D)_{ikjl}=\frac12\sum_s u_s J_{s,ik}J_{s,jl},
\]
where $ij$ is the row index and $kl$ is the column index.
These tensorial definitions may be evaluated after pointwise
diagonalization. They annihilate metric trace. On the symmetric
trace-free tensors, the nine matrices $I_iJ_j$ form an orthogonal
basis, with $T_D$ and $U_D$ eigenvalues $t_i$ and $u_j$.
For example, $T_D$ acts by
$-\frac12\sum_s t_s I_s(\,\cdot\,)I_s$; the quaternion relations
give the eigenvalue $t_i$ on $I_iJ_j$.
Thus
\[
 T_D+U_D\le 2(x+y)\operatorname{Id}<4\operatorname{Id}
\]
on this space.

The integrated curvature contraction has a useful sign and a precise
coefficient. Since $T_D$ annihilates trace, replace $S$ by
$\operatorname{Hess}\phi$ in its quadratic form. Integrate the
derivative $k$, use $\nabla_k(T_D)_{ikjl}=0$, and use skewness in
$i,k$ to obtain
\[
 \begin{aligned}
 \mathbf E\langle T_DS,S\rangle
 &=-\mathbf E (T_D)_{ikjl}\phi_l\nabla_k\nabla_i\phi_j\\
 &=\frac12\mathbf E (T_D)_{ikjl}R_{ikja}\phi_a\phi_l.
 \end{aligned}
\]
In this contraction the minus curvature block is orthogonal to the
plus block. Substituting their expansions gives
\[
 \frac12\sum_{ikj}(T_D)_{ikjl}R_{ikja}
 =\frac12\sum_s t_s(1+t_s)
                 \sum_j I_{s,jl}I_{s,ja}
 =3v^2\delta_{la}.
\]
The final equality uses $\sum t_s=0$ and $\sum t_s^2=6v^2$.
Consequently
\[
 \mathbf E\langle(T_D+U_D)S,S\rangle
 =3\mathbf E\bigl[(v^2+b^2)|\nabla\phi|^2\bigr].
\]
The pointwise upper operator bound, inserted into the positive
$|S|^2$ term of the Bochner identity, now yields
\[
 \mathbf E\bigl[(\Delta\phi)^2\bigr]\ge\mathbf E K|\nabla\phi|^2.
\]

For smooth $a$, the closed-manifold Poisson statement in
Appendix~\ref{app:analytic-preliminaries} gives a unique smooth
zero-mean solution of $\Delta\phi=a-\mathbf E a$. Integration
by parts and Cauchy--Schwarz with the pointwise weight $K$ give
\[
 \begin{aligned}
 \operatorname{Var}(a)
 &=-\mathbf E\langle\nabla a,\nabla\phi\rangle\\
 &\le
 \left(\mathbf E\frac{|\nabla a|^2}{K}\right)^{1/2}
 \left(\mathbf E K|\nabla\phi|^2\right)^{1/2}\\
 &\le
 \left(\mathbf E\frac{|\nabla a|^2}{K}\right)^{1/2}
 \operatorname{Var}(a)^{1/2}.
 \end{aligned}
\]
Division proves the assertion when the variance is nonzero; the
zero case is immediate. Smooth functions are dense in $H^1$ by the
local mollification argument recalled in
Appendix~\ref{app:analytic-preliminaries}. Since $K$ is bounded above
and below by positive constants, this approximation also passes the
weighted energy to the limit and proves the assertion on $H^1$.
\end{proof}

\subsection{The sharp scalar inequality and the Weyl gap}

The following is the four-dimensional Einstein case of Ilias's sharp
positive-Ricci Sobolev inequality~\cite[Theorem~3]{Ilias}. Its proof in
the present normalization is given in Appendix~\ref{app:sharp-sobolev}.
The nonlinear Hessian argument
belongs to the Bochner approach to elliptic rigidity; compare the
systematic treatment in~\cite{BidautVeronVeron}.

\begin{proposition}[Normalized Sobolev inequality]
\label{prop:sharp-sobolev}
For every real $f\in H^1(M)$,
\[
 \bigl(\mathbf E|f|^4\bigr)^{1/2}
 \le\mathbf E\left(f^2+\frac12|\nabla f|^2\right).
\]
\end{proposition}

The next proposition is the Weyl-norm gap established in
Gursky--LeBrun~\cite[Theorem~1]{GurskyLeBrun1999}; we give the derivation in
our normalization, including the equality case and the possible
zeros of the Weyl block.

\begin{proposition}[Weyl-norm gap]
\label{prop:weyl-gap}
For either Weyl block,
\begin{equation}
 v\equiv0\qquad\text{or}\qquad \mathbf E v^2\ge1.
 \tag{5}\label{eq:5}
\end{equation}
If $T\not\equiv0$ and $\mathbf E v^2=1$, then
$\nabla T=0$ and $v=\sigma=1$. The corresponding statements hold
for $U$.
\end{proposition}

\begin{proof}
Put $A_\varepsilon=v^2+\varepsilon^2$ and
$f_\varepsilon=A_\varepsilon^{1/6}$. These are smooth functions.
The equation for $\Delta(v^2)$ and \eqref{eq:2} give
\[
 \begin{aligned}
 \Delta f_\varepsilon
 =A_\varepsilon^{-5/6}
 \left[
  2(v^2-\sigma v^3)+\frac13e
       -\frac59\frac{v^2}{A_\varepsilon}|p|^2
 \right]
 \ge 2A_\varepsilon^{-5/6}(v^2-v^3).
 \end{aligned}
\]
Indeed, $\sigma\le1$, $v^2/A_\varepsilon\le1$, and
$e\ge5|p|^2/3$. Multiply by $f_\varepsilon$ and integrate:
\[
 \mathbf E|\nabla f_\varepsilon|^2
 \le 2\mathbf E\frac{v^3-v^2}{A_\varepsilon^{2/3}}
 \le 2\mathbf E v^{5/3}.
\]
Their $L^2$ norms are uniformly bounded as well.
Thus $f_\varepsilon$ is bounded in $H^1$, and its uniform limit
$f=v^{1/3}$ belongs to $H^1$ by weak compactness. This bound is
obtained directly from the regularized equation; it requires no
positive lower bound on $v$.

Apply Proposition~\ref{prop:sharp-sobolev} to $f_\varepsilon$
and use the preceding energy bound:
\[
 \begin{aligned}
 \bigl(\mathbf E f_\varepsilon^4\bigr)^{1/2}
 &\le \mathbf E\left[
       A_\varepsilon^{1/3}
       +\frac{v^3-v^2}{A_\varepsilon^{2/3}}\right]\\
 &=\mathbf E\frac{\varepsilon^2+v^3}{A_\varepsilon^{2/3}}.
 \end{aligned}
\]
The term $\varepsilon^2/A_\varepsilon^{2/3}$ is at most
$\varepsilon^{2/3}$, while
$v^3/A_\varepsilon^{2/3}\le v^{5/3}$.
Dominated convergence therefore gives
\[
 \bigl(\mathbf E f^4\bigr)^{1/2}
 \le\mathbf E(f^2v)
 \le\bigl(\mathbf E f^4\bigr)^{1/2}
       \bigl(\mathbf E v^2\bigr)^{1/2}.
\]
When $T\not\equiv0$, the common factor is positive, proving
\eqref{eq:5}.

If $\mathbf E v^2=1$, equality holds in Cauchy--Schwarz.
Thus $v=\alpha f^2$ almost everywhere for a constant $\alpha>0$.
It follows that $v$ takes only the values $0$ and $\alpha^3$
almost everywhere, and hence everywhere by continuity.
Connectedness and $T\not\equiv0$ force the positive constant value
on all of $M$. The moment condition gives $v=1$.
Equation~\eqref{eq:3} then becomes
$0=6(1-\sigma)+e$. Both terms are nonnegative, so
$\sigma=1$ and $e=0$. Smoothness gives $\nabla T=0$ everywhere.
\end{proof}

\subsection{The Euler and signature formulas}

\begin{proposition}[Characteristic identities]
\label{prop:characteristic-identities}
Let $n_+$ and $n_-$ be the finite nonnegative integer dimensions of
the spaces of harmonic self-dual and anti-self-dual two-forms. The Euler
characteristic and signature of $M$ are, respectively,
\[
 \chi=2+n_++n_-,\qquad \tau=n_+-n_-.
\]
Then
\begin{equation}
 \begin{aligned}
 8\pi^2\chi
   &=6\operatorname{Vol}(M)
        \bigl(1+\mathbf E(v^2+b^2)\bigr),\\
 12\pi^2\tau
   &=6\operatorname{Vol}(M)\,\mathbf E(v^2-b^2).
 \end{aligned}
 \tag{6}\label{eq:6}
\end{equation}
If $U=0$, then $n_-=0$; likewise $T=0$ implies $n_+=0$.
\end{proposition}

These are the Euler and signature formulas in the present curvature
normalization. Appendix~\ref{app:characteristic-identities} identifies
the topological integers through Hodge theory and proves the formulas
by the form-complex heat-supertrace method of
McKean--Singer~\cite{McKeanSinger1967}; compare the general formulations
and local coefficients in Gilkey~\cite[Theorems~1.7.6 and
4.8.16]{Gilkey1984}. That appendix also proves the zero-block kernel
vanishing statements, which will be needed after one-block vanishing.
